The analysis thus far reveals that individual designer responsibility is necessary
but insufficient. Designers should advocate for sustainable choices within their
professional sphere, but systemic change requires coordinated action at multiple
levels. This section first identifies individual actions designers can take
immediately, then proposes regulatory interventions necessary to make those
actions professionally viable. The things Designers can do can be split into two layers. First the individual and professional actions:
\begin{itemize}
  \item \textbf{Awareness:} You can't make ethical choices if you don't know the impact. Learn about the energy costs of the tools and patterns you use, teach others and support regulation.
\item \textbf{Digital Restraint:} Just because Midjourney can generate 50 logo options doesn't mean you should request 50. Ask yourself: 'Do I need this, or is it just easy to do?'
  \item \textbf{Green UX or Planet Centered Design}
\item \textbf{Reframe Sustainability as Performance}
\end{itemize}
In 2024 Kiourtis et al proposed best practices towards environmentally sustainable UX/UI and they showed that when applying these proposed best practices per view carbon footprint could be significantly reduced \cite{Kiourtis}. Similarily to this other Green UX or Planet Centered Design frameworks exist that should be known about. Even though these individual actions are important
the analysis thus far reveals a structural problem: designers bear responsibility for environmental consequences but are systematically denied the information necessary to act on that responsibility and act within systems where performance clashes with sustainability. Hence, I propose the following regulations.
\begin{itemize}
\item \textbf{Mandatory Disclosure Requirements:} Just as the European Union requires nutritional labeling on food products and energy ratings on appliances, digital services should display environmental impact metrics. Specifically, companies, applications and features should provide: (1) energy consumption per hour of use or per interaction (2) carbon emissions based on grid location where computable and (3) comparative information enabling informed choice.
\item \textbf{Default Setting Standards:} Based on the power default settings hold, requiring eco friendly default settings can have a big impact.
\end{itemize}
Users can make informed consumption decisions, exercising their autonomy. Designers can make informed implementation decisions, fulfilling professional ethics obligations currently impossible without data. Regulators can establish appropriate standards based on empirical baselines rather than contested estimates. Precedent exists: the EU AI Act \cite{EUAIAct2024} mandates transparency in agentic AI systems that involve critical infrastructure; Carbon labeling represents a logical extension of existing digital transparency requirements. Critics might object that measurements vary by device, network, and context, making standardized labels impractical. However, appliance energy labels similarly provide ranges. The goal is not false precision but directional transparency: making environmental trade-offs visible rather than invisible.
